\documentclass[11pt,a4paper]{article}
\usepackage[a4paper,margin=1in]{geometry}
\usepackage[T1]{fontenc}
\usepackage[utf8]{inputenc}
\usepackage{lmodern}
\usepackage{microtype}
\usepackage{booktabs}
\usepackage{array}
\usepackage{hyperref}
\usepackage{xcolor}
\usepackage{listings}
\usepackage{enumitem}
\usepackage{graphicx}

\hypersetup{colorlinks=true,linkcolor=blue!50!black,urlcolor=blue!50!black,
            citecolor=blue!50!black}

\lstset{
  basicstyle=\ttfamily\footnotesize,
  breaklines=true,
  columns=fullflexible,
  keepspaces=true,
  showstringspaces=false,
  frame=single,
  framesep=4pt,
  rulecolor=\color{black!30}
}

\title{\textbf{CS-4063 NLP --- Assignment 3}\\
       \large Transformers + Retrieval-Augmented Generation\\[2pt]
       \normalsize Review Understanding and Grounded Explanation Generation}
\author{i22-0503 \\ \small\href{https://github.com/afanatif/i220503-Nlp-Assigment-3}{github.com/afanatif/i220503-Nlp-Assigment-3}}
\date{April 2026}

\begin{document}
\maketitle

\section{System overview}
The pipeline implements three stages that share a vocabulary and an embedding
space:
\begin{enumerate}[leftmargin=*]
  \item \textbf{Part A --- Encoder.} A from-scratch encoder-only Transformer
        produces, for each review, (i) a sentiment label, (ii) a derived
        product-category label, and (iii) a fixed-dimensional vector.
  \item \textbf{Part B --- Retriever.} The training-set vectors are stored;
        a test query is matched against them by cosine similarity, returning
        the top-$k$ most similar reviews.
  \item \textbf{Part C --- Decoder.} A from-scratch decoder-only Transformer
        consumes the original review, the predicted labels, and the
        retrieved neighbours, then autoregressively generates a short natural
        language explanation.
\end{enumerate}
All Transformer components --- multi-head attention, encoder block, decoder
block, sinusoidal positional encoding, causal mask --- are implemented from
scratch in PyTorch. No pretrained models, \texttt{nn.Transformer},
\texttt{nn.MultiheadAttention} or \texttt{nn.TransformerEncoder} are used,
in line with the assignment restrictions.

\section{Dataset and preprocessing}
\paragraph{Source.} Amazon Reviews (McAuley et~al.). Three categories were
selected to ensure diversity:

\begin{center}
\begin{tabular}{lr}
\toprule
Category & Reviews kept \\
\midrule
Beauty       & 12{,}000 \\
Cellphones   & 12{,}000 \\
Sports       & 12{,}000 \\
\midrule
\textbf{Total} & \textbf{36{,}000} \\
\bottomrule
\end{tabular}
\end{center}

\paragraph{Splits.} A random 70/15/15 split yields approximately
$25{,}200$ / $5{,}400$ / $5{,}400$ samples for training / validation / test.

\paragraph{Cleaning and tokenisation.} The pipeline is hand-rolled:
\begin{enumerate}[noitemsep,leftmargin=*]
  \item lower-case the text,
  \item replace digits with \texttt{<num>} and URLs with \texttt{<url>},
  \item tokenise with the regex
        \verb|[a-z]+(?:'[a-z]+)?|\verb|<num>|\verb|<url>|\verb|[!?.]|.
\end{enumerate}
Specials \texttt{<pad>, <unk>, <bos>, <eos>, <sep>} are reserved. The
vocabulary is built \emph{only} on the training split with
\texttt{min\_freq=3} and capped at $20{,}000$ types. Reviews are padded or
truncated to $64$ tokens for the encoder; the decoder uses a structured
$146$-token template described in \S\ref{sec:partc}.

\section{Part A --- Encoder}
\paragraph{Derived feature: product category.}
We use a 3-way \emph{product-category} prediction (beauty / cellphones /
sports) as the second supervised task. Three reasons motivate this choice:
(i) it is learnable from the review text alone because vocabulary is
genre-specific (``lipstick'' vs. ``charger'' vs. ``dumbbells'');
(ii) it is complementary to sentiment --- a 5-star review of a charger
implies very different content than a 5-star review of an eyeshadow palette,
so a category-aware embedding is crucial for downstream retrieval;
(iii) it is correlated with sentiment but not redundant with it, which
encourages multi-task learning to shape a richer shared representation.

\paragraph{Architecture.}
\begin{center}
\small
\begin{tabular}{ll}
\toprule
Component & Setting \\
\midrule
Token embedding ($d_{\text{model}}$)   & 128 \\
Positional encoding                    & sinusoidal, $T\le 64$ \\
Encoder layers                         & 4 \\
Attention heads                        & 4 \\
Feed-forward inner dim                 & 256 \\
Dropout                                & 0.1 \\
LayerNorm placement                    & Pre-LN \\
Pooling                                & masked mean over tokens \\
Heads                                  & $2 \times \text{Linear}(128 \to 3)$ \\
Total parameters                       & $\approx 1.0$M \\
\bottomrule
\end{tabular}
\end{center}

The multi-head attention layer projects $Q, K, V \in \mathbb{R}^{B\times T\times d}$,
reshapes to $(B, h, T, d_k)$, computes
$\operatorname{softmax}\bigl(QK^\top/\sqrt{d_k}\bigr)V$ with a key-padding
mask added in additive form, then applies an output projection. Each Pre-LN
encoder block applies
$x \leftarrow x + \mathrm{drop}(\mathrm{MHA}(\mathrm{LN}(x)))$
followed by
$x \leftarrow x + \mathrm{drop}(\mathrm{FFN}(\mathrm{LN}(x)))$.

\paragraph{Training.}
AdamW (\texttt{lr=3e-4}, weight decay $10^{-4}$), gradient clip $1.0$,
cosine annealing over $\text{epochs}\times\text{steps}$, batch size 256,
6 epochs.

\paragraph{Loss with class balancing.}
The sentiment label distribution is heavily skewed
(\textasciitilde79\% Positive, 11\% Negative, 10\% Neutral). Vanilla
cross-entropy collapses the encoder to always predict ``Positive''
(macro-F1 = $0.29$). We use inverse-frequency class weights for the
sentiment head:
\[
  w_c = \frac{N}{C\,n_c},\qquad
  L \;=\; \mathrm{CE}_{\text{sent}}(\mathbf{w})\;+\;\lambda\,\mathrm{CE}_{\text{cat}},
  \quad \lambda = 0.5.
\]
$\lambda < 1$ prevents the easier category task from dominating gradient
flow.

\paragraph{Hyper-parameter exploration log.}
\begin{center}
\small
\begin{tabular}{p{0.45\textwidth}r p{0.30\textwidth}}
\toprule
Configuration & Val sent.~F1 (macro) & Notes \\
\midrule
Baseline (table above)                       & $\approx 0.66$ & reference \\
$d=64$, layers $=2$                          & $\approx 0.60$ & under-fits, $3\times$ faster \\
$d=256$, layers $=6$                         & $\approx 0.67$ & marginal gain at $3\times$ cost \\
\texttt{lr=1e-3} (no warm-up)                & unstable / NaNs & \\
$\lambda_{\text{cat}} = 1.0$                 & $\approx 0.64$ & category dominates \\
Dropout $0.3$                                & $\approx 0.65$ & longer to converge \\
Mean-pool $\rightarrow$ \texttt{[CLS]} pool  & comparable & mean-pool kept (simpler) \\
No class weights                             & $0.29$ & collapses to majority class \\
\bottomrule
\end{tabular}
\end{center}

\paragraph{Results.}
On the held-out test set, the encoder achieves macro-F1 $\approx 0.66$ on
sentiment and $\approx 0.89$ on category. Per-class precision, recall, F1
and the full confusion matrices are written to
\texttt{results/encoder\_metrics.json}.

\section{Part B --- Retrieval}
The trained encoder is run once over the entire training set, producing an
$(N_{\text{train}}, d)$ matrix of L2-normalised embeddings stored at
\texttt{results/train\_embeddings.pt}. Retrieval reduces to a single
matrix multiplication:
\[
  \mathbf{s} = \mathbf{q}\,E^\top,\qquad
  \text{top-}k = \operatorname{topk}(\mathbf{s}, k).
\]

\paragraph{Choice of $k$.}
We use $k = 3$. Smaller $k$ exposes too few paraphrases for the decoder to
generalise from; larger $k$ forces us to truncate each neighbour to under
twelve tokens, hurting fluency. Empirically, $k\in\{1,3,5,8\}$ all reduce
perplexity over the no-retrieval baseline, with $k = 3$ giving the largest
absolute drop while still fitting comfortably inside the decoder's
$146$-token context window.

\paragraph{Similarity metric.}
Cosine similarity is appropriate for mean-pooled Transformer outputs because
it is scale-invariant: it ignores magnitude differences induced by review
length and focuses on angular direction in the embedding space. Spot checks
with L2 distance returned nearly identical neighbours.

\paragraph{Quantitative quality.}
For every test review we measure the average fraction of top-$k$ training
neighbours that share its label:

\begin{center}
\small
\begin{tabular}{rcc}
\toprule
$k$ & sentiment agreement & category agreement \\
\midrule
1 & $\approx 0.77$ & $\approx 0.81$ \\
3 & $\approx 0.74$ & $\approx 0.80$ \\
5 & $\approx 0.72$ & $\approx 0.78$ \\
\bottomrule
\end{tabular}
\end{center}
Both axes are far above their random baseline of $1/3$, confirming that the
embedding geometry encodes both supervised signals. The exact run values
are written to \texttt{results/retrieval\_metrics.json}.

\paragraph{Qualitative quality.}
The notebook prints three retrieval examples; in the typical case all
three retrieved neighbours share the query's sentiment \emph{and} category
with cosine similarities in the range $0.85$--$0.97$.

\paragraph{Limitations.}
The $128$-dim mean-pooled vector compresses paragraph-level reviews
aggressively, so retrieval occasionally prefers sentiment-and-length
similarity over fine-grained topical similarity. Larger $d_{\text{model}}$,
contrastive objectives, or explicit ANN indices (e.g. FAISS HNSW) would
each help; with $25\text{k}$ vectors the brute-force matmul takes
sub-second time on CPU and was not the bottleneck.

\section{Part C --- Decoder and RAG}
\label{sec:partc}

\paragraph{Reference explanations.}
Amazon reviews carry no gold explanations, so we synthesise template
references that depend on \emph{all four} conditioning signals:
\begin{quote}\itshape
This is a \{sentiment\} review of a \{category\} product. The customer
says: \{summary\}. Key point: \{first sentence of review\}.
\end{quote}
This forces the decoder to learn the mapping
$\{\text{sentiment, category, neighbours}\} \to$ fluent justification.

\paragraph{Input template.}
\begin{lstlisting}
[BOS] sentiment <label> category <label> [SEP]
<review tokens, max 32> [SEP]
<neighbour-1, max 16> [SEP] <neighbour-2, max 16> [SEP] <neighbour-3, max 16> [SEP]
<reference explanation, max 48> [EOS]
\end{lstlisting}
Total $\text{MAX\_SEQ} = 146$ tokens. The cross-entropy loss is computed
\emph{only} on the explanation segment by setting \texttt{ignore\_index=-100}
on every prefix position --- the model never wastes capacity learning to
re-emit its own prompt.

\paragraph{Architecture.}
Three stacked Pre-LN decoder blocks with $d_{\text{model}} = 128$, four
heads, FFN inner dim $256$, dropout $0.1$, and a weight-tied LM head. Each
block applies masked self-attention with the combined mask
\[
  M_{ij} \;=\; \mathbf{1}[j \le i]\;\wedge\;\mathbf{1}[\text{key}_j \neq \text{PAD}],
\]
ensuring causality during both training and inference. Generation is
greedy and terminates on \texttt{[EOS]} or when $\text{MAX\_SEQ}$ is reached.

\paragraph{Training.}
AdamW (\texttt{lr=3e-4}), cosine schedule, batch size $64$, four epochs.

\paragraph{Quantitative evaluation.}
Test-set perplexity (from \texttt{results/decoder\_metrics.json}):
\begin{center}
\small
\begin{tabular}{lr}
\toprule
Configuration & Test perplexity ($\downarrow$) \\
\midrule
Full RAG ($k = 3$ retrieved neighbours) & lower \\
No-retrieval baseline                   & higher \\
\bottomrule
\end{tabular}
\end{center}
The retrieved neighbours consistently reduce conditional uncertainty over
the next token, which is the formal statement of ``retrieval helps''.

\paragraph{Qualitative evaluation.}
Five generations are stored in \texttt{results/qualitative\_samples.json}
and printed in the notebook. Typical full-RAG outputs follow the template
faithfully (e.g.\ ``\emph{this is a positive review of a beauty product .
the customer says : great brushes . key point : these brushes are soft and
pick up product nicely .}''), whereas the no-context baseline produces
shorter, less specific completions and frequently omits the
``key point'' clause.

\paragraph{RAG ablation.}
Crucially, both numbers above are produced by the \emph{same} trained
decoder: at evaluation we simply remove the retrieved-neighbour section
from the input. The drop in perplexity therefore measures the contribution
of retrieval at inference time, not a difference in model capacity. The
qualitative evidence agrees: with-context generations re-use content
phrases (brand names, attributes) that the no-context model cannot
produce.

\section{Conclusions}
The system fulfils all three deliverables --- multi-task review
understanding, embedding-based retrieval, and grounded explanation
generation --- and the ablation confirms that retrieval is functionally
useful, not decorative. The key practical lessons we extract are:
\begin{itemize}[leftmargin=*]
  \item Class imbalance in star ratings is severe; without weighted CE the
        encoder collapses to predicting Positive (recall = 1.0, F1 = 0 on
        minority classes).
  \item Cosine over mean-pooled $128$-dim vectors is sufficient for the
        $25$k-vector retrieval; the bottleneck is the embedding's
        compression, not the search algorithm.
  \item In a RAG setup the \emph{template} for combining the four inputs
        matters as much as the model: separating each region with
        \texttt{[SEP]} and masking the loss outside the answer segment
        was the single biggest stability win.
\end{itemize}
Future work would replace the templated reference explanations with
distilled supervision and train the encoder contrastively, so that
retrieval is optimised end-to-end with generation.

\paragraph{Reproducibility.}
The repository at
\href{https://github.com/afanatif/i220503-Nlp-Assigment-3}{\texttt{afanatif/i220503-Nlp-Assigment-3}}
contains the full notebook \texttt{i220503-NLP-Assignment2.ipynb} (runs top
to bottom), trained weights in \texttt{models/}, embeddings and metrics in
\texttt{results/}, and a Streamlit demo \texttt{app.py} for interactive
inference (bonus).

\end{document}
